\section{图与概率}

\subsection{图的记号与术语}

一个\emph{图}由顶点（或节点）集合 $V$ 与连接某些顶点对的边（或连线）集合 $E$ 构成。在我们的图中，顶点对应于变量（故通用符号为 $V$），边则表示在成对变量之间成立的某种关系，其解释随应用而异。由一条边相连的两个变量称为\emph{邻接}。

图中的每条边或者\emph{有向}（边上标有一个箭头），或者\emph{无向}（不带标记的连线）。在某些应用中，我们还会用“双向”边表示存在未被观测的共同原因（有时称为\textbf{混杂因子}）。这种边画成两端带箭头的虚线弧（见图 1.1(a)）。若所有边都是有向的（见图 1.1(b)），则得到\emph{有向图}。若把图 $G$ 中所有边上的箭头去掉，所得无向图称为 $G$ 的\emph{骨架}。图中的\emph{路径}是一列边（例如图 1.1(a) 中的 $((W,Z),(Z,Y),(Y,X),(X,Z))$），其中每条边的起点正是前一条边的终点。换言之，路径就是沿图中各边走出的一条不间断、不自交的路线，既可以顺着箭头走，也可以逆着箭头走。若路径中每条边都是从第一个顶点指向第二个顶点的箭头，则得到\emph{有向路径}。例如在图 1.1(a) 中，路径 $((W,Z),(Z,Y))$ 是有向的，而路径 $((W,Z),(Z,Y),(Y,X))$ 与 $((W,Z),(Z,X))$ 则不是。若图中两个顶点之间存在一条路径，则称这两个顶点\emph{连通}，否则称\emph{不连通}。

有向图可以包含\emph{有向环}（例如 $X\to Y,\ Y\to X$，表示互为因果或反馈过程），但不含自环（如 $X\to X$）。不含任何有向环的图（如图 1.1 中的两个图）称为\emph{无环的}。既是有向又是无环的图（图 1.1(b)）称为\emph{有向无环图}（directed acyclic graph，DAG），这类图将占据我们关于因果性的讨论的主要篇幅。我们自由地借用亲属关系的术语（如父节点、子节点、后代、祖先、配偶）来表示图中的各种关系。这些亲属关系沿图中的完整箭头定义，包括构成有向环的箭头，但忽略双向边与无向边。例如在图 1.1(a) 中，$Y$ 有两个父节点（$X$ 与 $Z$）、三个祖先（$X$、$Z$、$W$），没有子节点；而 $X$ 没有父节点（因而没有祖先），有一个配偶（$Z$），有一个子节点（$Y$）。图中的\emph{家族}是这样一个节点集，它包含某节点及其全部父节点。例如，$\{W\}$、$\{Z,W\}$、$\{X\}$、$\{Y,Z,X\}$ 就是图 1.1(a) 中的各个家族。

有向图中的节点若无父节点，称为\emph{根}；若无子节点，称为\emph{汇}。每个 DAG 至少有一个根、至少有一个汇。若一个连通 DAG 中每个节点至多有一个父节点，则称为\emph{树}；若树中每个节点至多有一个子节点，则称为\emph{链}。若图中每对节点都由一条边相连，则称为\emph{完全图}。例如图 1.1(a) 中的图是连通的但不完全，因为 $(W,X)$ 与 $(W,Y)$ 这两对节点并不邻接。

\begin{figure}[htbp]
\centering
\begin{tikzpicture}[causal,scale=0.95,every node/.style={transform shape}]
  % ---- (a) ----
  \begin{scope}
    \node[vtx] (W) at (0,2.2) {$W$};
    \node[vtx] (Z) at (-1.0,0.75) {$Z$};
    \node[vtx] (X) at (1.0,0.75) {$X$};
    \node[vtx] (Y) at (0,-0.9) {$Y$};
    \draw[ed] (W) -- (Z);
    \draw[ed] (W) -- (X);
    \draw[ed] (Z) -- (Y);
    \draw[ed] (X) -- (Y);
    \draw[bid,dashed] (Z) to[bend left=38] (X);
    \node[font=\small] at (0,-1.75) {(a)};
  \end{scope}
  % ---- (b) ----
  \begin{scope}[xshift=5.2cm]
    \node[vtx] (W2) at (0,2.2) {$W$};
    \node[vtx] (Z2) at (-1.0,0.75) {$Z$};
    \node[vtx] (X2) at (1.0,0.75) {$X$};
    \node[vtx] (Y2) at (0,-0.9) {$Y$};
    \draw[ed] (W2) -- (Z2);
    \draw[ed] (W2) -- (X2);
    \draw[ed] (Z2) -- (X2);
    \draw[ed] (Z2) -- (Y2);
    \draw[ed] (X2) -- (Y2);
    \node[font=\small] at (0,-1.75) {(b)};
  \end{scope}
\end{tikzpicture}
\caption{(a) 同时含有有向边与双向边的图。(b) 与 (a) 具有相同骨架的有向无环图（DAG）。}
\label{fig:1.1}
\end{figure}

\subsection{贝叶斯网络}

图在概率建模与统计建模中的作用有三：
\begin{enumerate}[label=\arabic*.]
\item 为表达实质性假设提供便利手段；
\item 便于经济地表示联合概率函数；
\item 便于从观测中进行高效推断。
\end{enumerate}
我们先从第 2 项谈起。

考虑为 $n$ 个二值变量指定一个任意联合分布 $P(x_1,\dots,x_n)$ 的任务。若要把 $P(x_1,\dots,x_n)$ 显式存储起来，就需要一张有 $2^n$ 个表目的表，无论按什么标准这都是难以想象的大数。当每个变量只依赖于少数几个其他变量时，就可以实现可观的节省。这类依赖信息使我们能够把大型分布函数分解为若干个小分布——每个只涉及少量变量——再把它们协调地拼装起来，以回答全局性的问题。图在这种分解中起着关键作用，因为它们直观地表示了在任一给定知识状态下彼此相关的那些变量集合。

研究者既用有向图也用无向图来促成这种分解。无向图有时称为\emph{马尔可夫网络}（Pearl 1988b），主要用于表示对称的空间关系（Isham 1981; Cox and Wermuth 1996; Lauritzen 1996）。有向图、尤其 DAG，则被用来表示因果或时间关系（Lauritzen 1982; Wermuth and Lauritzen 1983; Kiiveri et al. 1984），并被称为\textbf{贝叶斯网络}——这一术语由 Pearl (1985) 创制，用以强调三个方面：(1) 输入信息的主观性；(2) 以贝叶斯条件化作为更新信息的基础；(3) 因果推理模式与证据推理模式之别，这一区分正是托马斯·贝叶斯 1763 年论文的要旨。混合图（同时含有有向边与无向边）也被提出用于统计建模（Wermuth and Lauritzen 1990），但本书主要关注有向无环图，偶尔使用有向循环图表示反馈环。

有向无环图所提供的基本分解方案可说明如下。设有定义在 $n$ 个离散变量上的分布 $P$，我们可以把这些变量任意排序为 $X_1,X_2,\dots,X_n$。概率演算的链式法则（式 \textup{(1.13)}）总允许我们把 $P$ 分解为 $n$ 个条件分布之积：
\begin{equation}
P(x_1,\dots,x_n) = \prod_{j} P(x_j \mid x_1,\dots,x_{j-1}). \tag{1.30}
\end{equation}
现在假设某个变量 $X_j$ 的条件概率并不敏感于它的全部前驱，而只敏感于其中一小部分。换言之，假设一旦已知一组称为 $\PA_j$ 的前驱的取值，$X_j$ 就与其余一切前驱独立。于是我们可以在 \textup{(1.30)} 的乘积中写
\begin{equation}
P(x_j \mid x_1,\dots,x_{j-1}) = P(x_j \mid \pa_j), \tag{1.31}
\end{equation}
从而大大简化所需的输入信息。我们不必再对 $X_j$ 在其前驱 $X_1,\dots,X_{j-1}$ 的一切可能实现上指定条件概率，而只需考虑集合 $\PA_j$ 的各种可能实现。集合 $\PA_j$ 称为 $X_j$ 的\emph{马尔可夫父节点}，或简称\emph{父节点}。当我们围绕这一概念构造图时，这个名称的由来就变得清楚了。

\begin{definition}[马尔可夫父节点]
设 $V=\{X_1,\dots,X_n\}$ 是一个有序变量集，$P(v)$ 是这些变量上的联合概率分布。若 $\PA_j$ 是 $X_j$ 的一个极小前驱集合，使得在给定 $\PA_j$ 的条件下 $X_j$ 与其余一切前驱独立，则称 $\PA_j$ 为 $X_j$ 的\emph{马尔可夫父节点}。换言之，$\PA_j$ 是 $\{X_1,\dots,X_{j-1}\}$ 的任一满足
\begin{equation}
P(x_j \mid \pa_j) = P(x_j \mid x_1,\dots,x_{j-1}) \tag{1.32}
\end{equation}
的子集，且 $\PA_j$ 的任何真子集都不满足 \textup{(1.32)}。\footnote{小写符号（如 $x_j,\pa_j$）表示相应变量（如 $X_j,\PA_j$）的某个具体实现。}
\end{definition}

定义 1.2.1 为每个变量 $X_j$ 指定了一个精选的前驱集 $\PA_j$，它足以确定 $X_j$ 的概率；一旦知道了父节点集 $\PA_j$ 的取值 $\pa_j$，再知道其他前驱的取值便是多余的。这一指定可以用 DAG 的形式表示：变量用节点表示，并从父节点集 $\PA_j$ 的每个节点向子节点 $X_j$ 画箭头。定义 1.2.1 还提示了一种构造此类 DAG 的简单递归方法：先从数对 $(X_1,X_2)$ 开始，当且仅当两变量相依时从 $X_1$ 向 $X_2$ 画箭头。接着处理 $X_3$：若 $X_3$ 与 $\{X_1,X_2\}$ 独立，则不画箭头；否则考察是 $X_2$ 把 $X_3$ 与 $X_1$ 屏蔽开，还是 $X_1$ 把 $X_3$ 与 $X_2$ 屏蔽开。前一种情形从 $X_2$ 向 $X_3$ 画箭头，后一种情形从 $X_1$ 向 $X_3$ 画箭头。若找不到任何屏蔽条件，则从 $X_1$ 与 $X_2$ 同时向 $X_3$ 画箭头。一般地：在构造的第 $j$ 步，我们选取 $X_j$ 的任一极小前驱集，使该集合把 $X_j$ 与其其余前驱屏蔽开（如式 \textup{(1.32)}），称此集为 $\PA_j$，并从 $\PA_j$ 的每个成员向 $X_j$ 画箭头。所得的有向无环图称为\emph{贝叶斯网络}，其中从 $X_i$ 指向 $X_j$ 的箭头把 $X_i$ 指派为 $X_j$ 的马尔可夫父节点，这与定义 1.2.1 一致。

可以证明（Pearl 1988b），只要分布 $P(v)$ 严格正（即不含逻辑约束或定义性约束），父节点集 $\PA_j$ 就是唯一的；所谓严格正，是指变量的每一种配置 $v$ 无论多么不可能，都具有某个有限的概率。在此条件下，与 $P(v)$ 相联系的贝叶斯网络在给定变量次序后便是唯一的。

图 1.2 展示了一个简单而典型的贝叶斯网络。它描述了一年中的季节（$X_1$）、是否下雨（$X_2$）、洒水器是否打开（$X_3$）、路面是否会湿（$X_4$）以及路面是否会打滑（$X_5$）之间的关系。图中除根变量 $X_1$ 可取春、夏、秋、冬四个值之外，其余变量都是二值的（取真或假）。该网络是依据定义 1.2.1、以因果直觉为指引构造出来的。例如，$X_1$ 与 $X_5$ 之间没有直接连线，这刻画了我们的理解：季节变化对路面打滑的影响是以其他条件（如路面湿度）为中介的。这一直觉与 \textup{(1.32)} 的独立条件相吻合，因为一旦知道 $X_4$，$X_5$ 便与 $\{X_1,X_2,X_3\}$ 独立。

\begin{figure}[htbp]
\centering
\begin{tikzpicture}[causal,scale=1.0,every node/.style={transform shape}]
  \node[vtx] (X1) at (0,0) {$X_1$};  \node[right=1.5mm of X1,font=\small] {SEASON};
  \node[vtx] (X3) at (-2.0,-1.35) {$X_3$}; \node[left=1.5mm of X3,font=\small,align=right] {SPRINKLER};
  \node[vtx] (X2) at (2.0,-1.35) {$X_2$};  \node[right=1.5mm of X2,font=\small] {RAIN};
  \node[vtx] (X4) at (0,-2.7) {$X_4$};  \node[right=1.5mm of X4,font=\small] {WET};
  \node[vtx] (X5) at (0,-4.05) {$X_5$}; \node[right=1.5mm of X5,font=\small] {SLIPPERY};
  \draw[ed] (X1) -- (X3);
  \draw[ed] (X1) -- (X2);
  \draw[ed] (X3) -- (X4);
  \draw[ed] (X2) -- (X4);
  \draw[ed] (X4) -- (X5);
\end{tikzpicture}
\caption{表示五个变量之间依赖关系的贝叶斯网络。}
\label{fig:1.2}
\end{figure}

由定义 1.2.1 所蕴含的构造，把贝叶斯网络定义为沿构造次序承载条件独立关系的载体。显然，每个满足 \textup{(1.32)} 的分布必定（利用 \textup{(1.30)} 的链式法则）分解为乘积
\begin{equation}
P(x_1,\dots,x_n) = \prod_{i} P(x_i \mid \pa_i). \tag{1.33}
\end{equation}
例如，图 1.2 中的 DAG 诱导出分解
\begin{equation}
P(x_1,x_2,x_3,x_4,x_5) = P(x_1)P(x_2\mid x_1)P(x_3\mid x_1)P(x_4\mid x_2,x_3)P(x_5\mid x_4). \tag{1.34}
\end{equation}
\textup{(1.33)} 中的乘积分解不再与次序有关：给定 $P$ 与 $G$，我们可以检验 $P$ 是否分解为 \textup{(1.33)} 所给的乘积，而无需援引任何变量次序。因此可以断定：DAG $G$ 成为概率分布 $P$ 的贝叶斯网络的一个必要条件，是 $P$ 容许 $G$ 所支配的、如 \textup{(1.33)} 那样的乘积分解。

\begin{definition}[马尔可夫相容]
若概率函数 $P$ 相对于 DAG $G$ 容许 \textup{(1.33)} 的分解，则称 $G$ \emph{表示} $P$，称 $G$ 与 $P$ \emph{相容}，或称 $P$ 相对于 $G$ 是\emph{马尔可夫}的。\footnote{后一种说法在晚近文献中似更流行（如 Spirtes et al. 1993; Lauritzen 1996）。Pearl (1988b, p.~116) 曾用“$G$ 是 $P$ 的 I-图（I-map）”。}
\end{definition}

确定 DAG 与概率之间的相容性在统计建模中之所以重要，主要因为相容性是 DAG $G$ 能够解释由 $P$ 所表示的一批经验数据的充要条件，也就是说，是存在一个能够生成 $P$ 的随机过程的充要条件（如 Pearl 1988b, pp.~210--23）。若每个变量 $X_i$ 的取值都按某个概率 $P_i(x_i\mid\pa_i)$ 随机选取，且仅依据此前已选定的 $\PA_i$ 之值 $\pa_i$，则所生成实例 $x_1,x_2,\dots,x_n$ 的总体分布 $P$ 将相对于 $G$ 是马尔可夫的。反之，若 $P$ 相对于 $G$ 是马尔可夫的，则存在一组概率 $P_i(x_i\mid\pa_i)$，按之选取各变量 $X_i$ 的取值，可使所生成实例 $x_1,x_2,\dots,x_n$ 的分布恰好等于 $P$。（事实上，$P_i(x_i\mid\pa_i)$ 的正确取法就是 $P(x_i\mid\pa_i)$。）

刻画与 DAG $G$ 相容的分布集的一种便利方式，是列出每个这样的分布都必须满足的（条件）独立关系集。这些独立关系可以借助一个称为\emph{d-分离}的图形准则从 DAG 中读出（Pearl 1988b；字母 d 表示方向性 directional），它将在本书的许多讨论中扮演重要角色。

\subsection{d-分离准则}

考虑三个互不相交的变量集合 $X$、$Y$、$Z$，它们在有向无环图 $G$ 中表示为节点。要检验在任何与 $G$ 相容的分布中 $X$ 是否在给定 $Z$ 时与 $Y$ 独立，我们需要检验对应于变量 $Z$ 的节点是否“阻断”了从 $X$ 中节点到 $Y$ 中节点的一切路径。所谓\emph{路径}，指图中一列连续的边（方向任意）；所谓\emph{阻断}，则应理解为切断由这些路径所连接的变量之间的信息流（或依赖流），其定义如下。

\begin{definition}[d-分离]
一条路径 $p$ 被节点集 $Z$ \emph{d-分离}（或\emph{阻断}），当且仅当
\begin{enumerate}[label=\arabic*.]
\item $p$ 含有一条链 $i\to m\to j$ 或一个分叉 $i\leftarrow m\to j$，使得中间节点 $m$ 属于 $Z$；或者
\item $p$ 含有一个倒分叉（或碰撞点）$i\to m\leftarrow j$，使得中间节点 $m$ 不属于 $Z$，且 $m$ 的任何后代都不属于 $Z$。
\end{enumerate}
若 $Z$ 阻断了从 $X$ 中任一节点到 $Y$ 中任一节点的一切路径，则称集合 $Z$ \emph{d-分离} $X$ 与 $Y$。
\end{definition}

d-分离背后的直觉很简单；若赋予图中箭头以因果含义，就最容易领会。在因果链 $i\to m\to j$ 与因果分叉 $i\leftarrow m\to j$ 中，两个端点变量是边缘相依的，但一旦对中间变量取条件（即已知其取值），二者就彼此独立（即被阻断）。形象地说，对 $m$ 取条件似乎“阻断”了沿路径的信息流，因为已知 $m$ 之后，获知 $i$ 对 $j$ 的概率不再有影响。倒分叉 $i\to m\leftarrow j$ 表示两个原因具有共同结果，其作用恰好相反：若两个端点变量（边缘地）独立，则一旦对中间变量（即共同结果）或其任一后代取条件，它们就变得相依（即通过未阻断路径连通）。这一点可在图 1.2 的语境中加以确认。一旦知道季节，$X_3$ 与 $X_2$ 便独立（假定洒水器是依据季节事先设定的）；而一旦发现路面湿了或打滑，$X_2$ 与 $X_3$ 就变得相依，因为排除其中一种解释会提高另一种解释的概率。

在图 1.2 中，$X=\{X_2\}$ 与 $Y=\{X_3\}$ 被 $Z=\{X_1\}$ d-分离，因为连接 $X_2$ 与 $X_3$ 的两条路径都被 $Z$ 阻断。路径 $X_2\leftarrow X_1\to X_3$ 被阻断，因为它是分叉，其中间节点 $X_1$ 属于 $Z$；路径 $X_2\leftarrow X_4\to X_3$ 被阻断，因为它是倒分叉，其中间节点 $X_4$ 及其所有后代都在 $Z$ 之外。然而，$X$ 与 $Y$ 不被集合 $Z'=\{X_1,X_5\}$ d-分离：路径 $X_2\to X_4\leftarrow X_3$（一个倒分叉）不被 $Z'$ 阻断，因为 $X_5$ 是中间节点 $X_4$ 的后代，而 $X_5\in Z'$。打个比方，获知后果 $X_5$ 的取值会使它的两个原因 $X_2$ 与 $X_3$ 变得相依，仿佛沿汇聚于 $X_4$ 的各箭头打开了一条通路。

初看之下，读者也许会觉得有点奇怪：对一个并不位于被阻断路径上的节点取条件，竟能解除该路径的阻断。然而这正对应着因果关系的一种普遍模式：对两个独立原因的某个共同后果进行观测，往往会使这两个原因变得相依，因为在该后果已经发生的条件下，关于其中一个原因的信息会使另一个原因变得更为可能或更不可能。这一模式在统计学文献中称为\emph{选择偏倚}或\emph{伯克森悖论}（Berkson 1946），在人工智能中称为\emph{消释效应}（explaining away，Kim and Pearl 1983）。例如，若某研究生院的录取标准是“本科成绩优异”或“具有特殊音乐才能”，那么在该校学生群体中，这两种属性会被发现呈（负）相关，即便它们在总体人群中并不相关。确实，成绩偏低的学生很可能在音乐上天赋异禀，这正好解释了他们的被录取。

图 1.3 展示了更为复杂的 d-分离例子：例 (a) 含有一条双向弧 $Z_1\leftrightarrow Z_3$，例 (b) 涉及一个有向环 $X\to Z_2\to Z_1\to X$。在图 1.3(a) 中，当 $\{Z_1,Z_2,Z_3\}$ 中无一被测量时，$X$ 与 $Y$ 之间的两条路径都被阻断。然而，一旦测量了 $Z_1$，路径 $X\to Z_1\leftrightarrow Z_3\leftarrow Y$ 就被解除阻断。这是因为 $Z_1$ 同时解除了 $Z_1$ 与 $Z_3$ 两处的“碰撞”：前者是因为 $Z_1$ 正是该碰撞点的汇合节点，后者是因为 $Z_1$ 经路径 $Z_1\leftarrow Z_2\leftarrow Z_3$ 成为碰撞节点 $Z_3$ 的后代。在图 1.3(b) 中，$X$ 与 $Y$ 不能被任何节点集 d-分离，包括空集。若对 $Z_2$ 取条件，我们阻断了路径 $X\leftarrow Z_1\leftarrow Z_2\leftarrow Y$，却解除了路径 $X\to Z_2\leftarrow Y$ 的阻断；若对 $Z_1$ 取条件，我们再次阻断了路径 $X\leftarrow Z_1\leftarrow Z_2\leftarrow Y$，同时解除了路径 $X\to Z_2\leftarrow Y$ 的阻断，因为 $Z_1$ 是碰撞节点 $Z_2$ 的后代。

\begin{figure}[htbp]
\centering
\begin{tikzpicture}[causal,scale=1.0,every node/.style={transform shape}]
  % ---- (a) ----
  \begin{scope}
    \node[circle,fill,inner sep=1.4pt] (X) at (0,0) {};   \node[font=\small,below=1.6mm] at (0,0) {$X$};
    \node[circle,fill,inner sep=1.4pt] (A1) at (2.0,0) {}; \node[font=\small,below=1.6mm] at (2.0,0) {$Z_1$};
    \node[circle,fill,inner sep=1.4pt] (A2) at (3.6,0) {}; \node[font=\small,below=1.6mm] at (3.6,0) {$Z_2$};
    \node[circle,fill,inner sep=1.4pt] (A3) at (5.2,0) {}; \node[font=\small,below=1.6mm] at (5.2,0) {$Z_3$};
    \node[circle,fill,inner sep=1.4pt] (Y) at (7.0,0) {};  \node[font=\small,below=1.6mm] at (7.0,0) {$Y$};
    \draw[ed] (X) -- (A1);
    \draw[ed] (A2) -- (A1);
    \draw[ed] (A3) -- (A2);
    \draw[ed] (Y) -- (A3);
    \draw[bid,dashed] (A1) to[bend left=40] (A3);
    \node[font=\small] at (3.5,-1.55) {(a)};
  \end{scope}
  % ---- (b) ----
  \begin{scope}[xshift=8.4cm]
    \node[circle,fill,inner sep=1.4pt] (Xb) at (0,0) {};
    \node[font=\small,below=1.5mm] at (0,0) {$X$};
    \node[circle,fill,inner sep=1.4pt] (Z1b) at (1.45,1.25) {};
    \node[font=\small,above=1.2mm] at (1.45,1.25) {$Z_1$};
    \node[circle,fill,inner sep=1.4pt] (Z2b) at (2.45,0) {};
    \node[font=\small,below=1.5mm] at (2.45,0) {$Z_2$};
    \node[circle,fill,inner sep=1.4pt] (Yb) at (4.0,0) {};
    \node[font=\small,below=1.5mm] at (4.0,0) {$Y$};
    \draw[ed] (Z1b) -- (Xb);
    \draw[ed] (Z2b) -- (Z1b);
    \draw[ed] (Xb) -- (Z2b);
    \draw[ed] (Yb) -- (Z2b);
    \node[font=\small] at (2.0,-1.35) {(b)};
  \end{scope}
\end{tikzpicture}
\caption{说明 d-分离的图。在 (a) 中，$X$ 与 $Y$ 在给定 $Z_2$ 时 d-分离，在给定 $Z_1$ 时 d-连通。在 (b) 中，$X$ 与 $Y$ 不能被任何节点集 d-分离。}
\label{fig:1.3}
\end{figure}

d-分离与条件独立之间的联系，由 Verma and Pearl (1988；另见 Geiger et al. 1990) 的下列定理建立。

\begin{theorem}[d-分离的概率含义]
若在 DAG $G$ 中集合 $X$ 与 $Y$ 被 $Z$ d-分离，则在与 $G$ 相容的每个分布中，$X$ 在给定 $Z$ 时与 $Y$ 独立。反之，若在 DAG $G$ 中 $X$ 与 $Y$ 不被 $Z$ d-分离，则至少存在一个与 $G$ 相容的分布，使 $X$ 与 $Y$ 在给定 $Z$ 时相依。
\end{theorem}

定理 1.2.4 的逆命题部分其实要强得多——d-分离的缺失意味着在与 $G$ 相容的\emph{几乎所有}分布中都有相依。原因在于，要使图中某条未被阻断的路径上出现独立，需要对参数作精确的调节，而这种调节在实践中不大可能出现（见 Spirtes et al. 1993 及 \S 2.4、\S 2.9.1）。

为区分条件独立这一概率概念与 d-分离这一图形概念，后者采用记号 $(X\indep Y\mid Z)_G$。于是定理 1.2.4 可更简洁地表述如下。

\begin{theorem}
对 DAG $G$ 中任意三个互不相交的节点子集 $(X,Y,Z)$ 以及一切概率函数 $P$，有：
\begin{enumerate}[label=(\roman*)]
\item 若 $G$ 与 $P$ 相容，则 $(X\indep Y\mid Z)_G\Rightarrow(X\indep Y\mid Z)_P$；
\item 若 $(X\indep Y\mid Z)_P$ 在与 $G$ 相容的一切分布中都成立，则 $(X\indep Y\mid Z)_G$。
\end{enumerate}
\end{theorem}

Lauritzen 等 (1990) 基于祖先图的概念提出了另一种 d-分离检验。为检验 $(X\indep Y\mid Z)_G$，从 $G$ 中删去除 $\{X,Y,Z\}$ 及其祖先之外的一切节点，把每对具有共同子节点的节点用一条边连接起来，并去掉所有弧上的箭头。于是 $(X\indep Y\mid Z)_G$ 成立，当且仅当在所得无向图中 $Z$ 截断了 $X$ 与 $Y$ 之间的一切路径。

注意，构造图时所依据的次序并不进入 d-分离准则；决定概率 $P$ 必须满足哪些独立关系的，只是所得图的拓扑结构。事实上，可以证明下述定理（Pearl 1988b, p.~120）。

\begin{theorem}[有序马尔可夫条件]
概率分布 $P$ 相对于 DAG $G$ 为马尔可夫的充要条件是：在某个与 $G$ 的箭头相一致的变量次序下，每个变量在给定其在 $G$ 中的父节点时，与它的所有前驱独立。
\end{theorem}

由该定理可得一个与次序无关的判据，用以确定给定概率 $P$ 是否相对于给定 DAG $G$ 为马尔可夫的。

\begin{theorem}[父节点马尔可夫条件]
概率分布 $P$ 相对于 DAG $G$ 为马尔可夫的充要条件是：每个变量在给定其父节点时，与它的所有非后代（在 $G$ 中）独立。（谈到 $X_i$ 的“非后代”时，我们把 $X_i$ 本身排除在外。）
\end{theorem}

这一条件被 Kiiveri 等 (1984) 与 Lauritzen (1996) 称为“局部”马尔可夫条件，有时也被直接取作贝叶斯网络的定义（Howard and Matheson 1981）。然而在实践中，有序马尔可夫条件更便于使用。

由 d-分离还可得到另一条重要性质：一个判据，用以确定两个给定 DAG 是否\emph{观测等价}——也就是说，是否每个与其中一个 DAG 相容的概率分布也都与另一个相容。

\begin{theorem}[观测等价]
两个 DAG 观测等价，当且仅当它们具有相同的骨架和相同的 v 结构集合；所谓 v 结构，是指两条收敛箭头，其箭尾之间没有箭头相连（Verma and Pearl 1990）。\footnote{同样的判据由 Frydenberg (1990) 在链图的语境中独立导出，那里假定了严格正性。}
\end{theorem}

观测等价为我们从纯概率信息推断方向性的能力划定了界限。两个观测等价的网络，若不动用操纵性实验或时间信息，就无法区分。例如，把图 1.2 中 $X_1$ 与 $X_2$ 之间箭头的方向反转，既不会引入也不会破坏 v 结构。因此这种反转得到的是观测等价的网络，连线 $X_1\to X_2$ 的方向性无法由概率信息确定。然而，箭头 $X_2\to X_4$ 与 $X_4\to X_5$ 的性质不同：无论如何反转其方向，都会产生新的 v 结构。于是我们看到，某些概率函数 $P$（例如图 1.2 中贝叶斯网络据以构造的那个分布）在不伴随时间信息时，能够约束图中某些箭头的方向性。这种方向性约束的确切含义——以及利用这些约束从数据推断因果关系的可能性——将在第二章形式化。

\subsection{贝叶斯网络中的推断}

贝叶斯网络发展于 20 世纪 80 年代初，用以促进人工智能（AI）系统中的预测与“回溯”任务。在这些任务中，必须为传入的观测找到一个自洽的解释，使之既与观测一致，又与手头的先验信息一致。在数学上，任务归结为计算 $P(y\mid x)$，其中 $X$ 是一组观测，$Y$ 是一组被认为对预测或诊断重要的变量。

给定联合分布 $P$，$P(y\mid x)$ 的计算在概念上是平凡的，直接应用贝叶斯法则即得
\begin{equation}
P(y\mid x) = \frac{\sum_{s} P(y,x,s)}{\sum_{y,s} P(y,x,s)}, \tag{1.35}
\end{equation}
其中 $S$ 表示除 $X$ 与 $Y$ 之外的一切变量。由于每个贝叶斯网络都定义了一个联合概率 $P$（由 \textup{(1.33)} 中的乘积给出），显然 $P(y\mid x)$ 可以由 DAG $G$ 及其各家族上定义的条件概率 $P(x_i\mid\pa_i)$ 算出。

然而，挑战在于高效地完成这些计算，并且要在网络拓扑所提供的表示层面上完成。对于需要为自身推理过程生成解释的系统而言，后一点尤为重要。尽管这类推断技术对我们的因果性讨论并非必需，我们仍将简要概述，因为它们展示了：(i) 以图的形式组织概率知识的有效性；(ii) 在这种组织上进行自洽的概率计算（及其近似）的可行性。细节可参看所引文献。

最早提出的贝叶斯网络概率计算算法采用消息传递结构，且仅限于树状网络（Pearl 1982; Kim and Pearl 1983）。在这一技术中，每个变量被赋予一个简单的处理器，并允许它异步地向邻居传递消息，直至达到平衡（在有限步之内）。此后发展出的方法把这种树传播（及其某些同步变体）推广到一般网络。其中最流行的是 Lauritzen and Spiegelhalter (1988) 的连接树传播法，以及割集条件化法（Pearl 1988b, pp.~204--10; Jensen 1996）。在连接树方法中，我们把网络分解为若干构成树结构的簇（如团），然后把每个簇中的变量集当作一个复合变量，使它能够向邻居（同样是复合变量）传递消息。例如，图 1.2 的网络可以组织成三个簇构成的马尔可夫相容链：
\[
\{X_1,X_2,X_3\}\ \to\ \{X_2,X_3,X_4\}\ \to\ \{X_4,X_5\}.
\]
在割集条件化方法中，我们对一组变量赋予具体取值（实例化），使剩余网络成为树状；然后在该树上进行传播，再选取新的实例化，直至穷尽所有实例化；最后对结果取平均。例如在图 1.2 中，若把 $X_1$ 实例化为任一具体值（比如 $X_1=\text{夏}$），就切断了 $X_2$ 与 $X_3$ 之间的通路，剩余网络便成为树状结构。割集条件化方法的主要优点在于其存储空间需求极小（与网络规模成线性关系），而连接树方法的需求可能是指数级的。也有人提出了这两种基本算法的混合组合（Shachter et al. 1994; Dechter 1996），以便在存储与时间之间灵活权衡（Darwiche 2009）。

一般网络中的推断是“NP-难的”（Cooper 1990），但上述每种方法的计算复杂度都可以在实际处理之前估计出来。当估计值超出合理界限时，可以改用随机模拟之类的近似方法（Pearl 1988b, pp.~210--23）。该方法利用网络的拓扑结构，在局部变量子集上执行吉布斯抽样，既可串行也可并行。

DAG 的更多性质及其在专家系统证据推理中的应用，见 Pearl (1988b)、Lauritzen and Spiegelhalter (1988)、Pearl (1993a)、Spiegelhalter 等 (1993)、Heckerman 等 (1995) 以及 Darwiche (2009)。
